\documentclass[12pt]{article}
\usepackage[left = 0.6in, right = 0.6in, top = 0.4in, bottom = 1.5in]{geometry} 
\usepackage{pgfplots, fontawesome5, booktabs, xcolor, hyperref, amsmath,amsthm,amssymb,amsfonts, enumitem, fancyhdr, color, comment, graphicx, environ, actuarialsymbol}
\usepackage[dvipsnames]{xcolor}
\pagestyle{fancy}
\setlength{\headheight}{65pt}
%%%%%%%%%%%%%%%%%%%%%%%%%%%%%%%%%%%%%%%%%%%%%
\lhead{}
\rhead{Quotient functions}
%%%%%%%%%%%%%%%%%%%%%%%%%%%%%%%%%%%%%%%%%%%%%
\newcommand{\emptygrid}[3]{ % #1: range, #2: y-label, #3: "labels" or "nolabels"
    \begin{tikzpicture}
    \begin{axis}[
        axis lines = middle,
        xlabel = $x$,
        ylabel = $#2$,
        xmin=-#1, xmax=#1,
        ymin=-#1, ymax=#1,
        grid = both,
        grid style = {line width=.1pt, draw=gray!30},
        xtick distance = 5,
        ytick distance = 5,
        minor tick num = 4,
        tick label style = {font=\tiny},
        width=8cm, height=8cm,
        axis equal image,
        % Logic to hide labels if #3 is "nolabels"
        xticklabels={ \ifx\empty#3\empty\else\ifnum\pdfstrcmp{#3}{nolabels}=0 \else \fi\fi },
        yticklabels={ \ifx\empty#3\empty\else\ifnum\pdfstrcmp{#3}{nolabels}=0 \else \fi\fi }
    ]
    \end{axis}
    \end{tikzpicture}
}
%%%%%%%%%%%%%%%%%%%%%%%%%%%%%%%%%%%%%%%%%%%%%
\newtheorem{theorem}{Theorem}[section]
\begin{document}

\section{Quotient/reciprocal functions}
Functions where there is another function in the denominator
eg:

\begin{align*}
f(x) = \frac{1}{2x^{2}-4x+3}
\end{align*}

Try plotting it:
\begin{center}
\emptygrid{10}{f}{labels}
\end{center}
Try plotting this by hand:
\begin{align*}
y = \frac{1}{x^{2}-4x + 4}
\end{align*}

    \begin{center}
\emptygrid{10}{f}{nolabels}
\end{center}
Strategy?
\begin{itemize}
    \item Plot denominator first
    \item Gauge whether the function shoots up to $+\infty$ or $-\infty$ based on whether the denominator is positive or negative
\end{itemize}

\subsection{Notable reciprocal functions}
Trigonometric reciprocal functions: (Plot each one, note where the asymptotes are)
\begin{itemize}
    \item $\frac{1}{\sin(x)} = \csc(x)$
    
    \item $\frac{1}{\cos(x)} = \sec(x)$

    \item $\frac{1}{\tan(x)} = \cot(x)$
\end{itemize}
(give them new names so as to not mistake them for inverses)
    \begin{center}
\emptygrid{10}{f}{nolabels} \emptygrid{10}{f}{nolabels}
\emptygrid{10}{f}{nolabels}
\end{center}

Transformations work the same as normal functions

\newpage
\subsection{Transformations on reciprocal trig functions}
Plot the following over $(-\pi ,2\pi )$
\begin{itemize}
    \item $f(x) = 2\csc(x+3)$

    \emptygrid{10}{f}{nolabels}
    \item $g(x) = -\frac{1}{2}\sec(2x- \frac{\pi}{4}) + 1$
    
    \emptygrid{10}{f}{nolabels}
\end{itemize}

\end{document}